\documentclass[twocolumn]{article}

% ---- Packages ----
\usepackage[utf8]{inputenc}
\usepackage[T1]{fontenc}
\usepackage{lmodern}
\usepackage{microtype}
\usepackage{graphicx}
\usepackage{booktabs}
\usepackage{hyperref}
\usepackage{xcolor}
\usepackage[margin=2cm]{geometry}
\usepackage{enumitem}
\usepackage{amsmath}
\usepackage{natbib}
\bibliographystyle{unsrtnat}

\hypersetup{
  colorlinks=true,
  linkcolor=blue!60!black,
  citecolor=blue!60!black,
  urlcolor=blue!60!black
}

% ---- Title ----
\title{BIME: A Browser-Based Molecular Editor with Built-in Substructure Intelligence}

\author{
  Syed Asad Rahman\thanks{Corresponding author: \texttt{s9asad@gmail.com}} \\
  BioInception PVT LTD, Cambridge, UK
}

\date{}

\begin{document}

\maketitle

% ====================================================================
\begin{abstract}
BIME (BioInception Molecular Editor) is a browser-based, SVG-rendered
molecular structure editor written in pure JavaScript with zero external
dependencies.  Unlike existing browser editors that focus on drawing
alone, BIME integrates a full cheminformatics engine---SMSD
Pro---directly in the client, providing maximum common substructure
(MCS) detection, substructure search via VF2++, molecular fingerprints
(path-based, ECFP, FCFP, topological torsion), and Tanimoto/Dice/cosine
similarity scoring without any server round-trip.  BIME implements the
IUPAC 2013 Cahn--Ingold--Prelog recommendations (Rules~1--4a) for
assigning R/S and E/Z stereodescriptors using a digraph-based algorithm
with phantom-node expansion.  A force-directed layout engine with
PrEd-inspired crossing preservation, SMACOF stress majorisation, and 45
scaffold templates produces publication-quality 2D coordinates.  The
editor supports SMILES, SMARTS, MOL V2000, and MOL V3000 formats, and
includes 552 built-in drug and biomolecule structures.  BIME is released
under the Apache License 2.0, requires no build tools or package
managers, and can be embedded in any web page with a single
\texttt{<script>} tag.  It is freely available at
\url{https://asad.github.io/bime-dist/}.
\end{abstract}

\noindent\textbf{Keywords:} molecular editor, maximum common
substructure, substructure search, molecular fingerprints,
Cahn--Ingold--Prelog, SMILES, JavaScript, SVG

% ====================================================================
\section{Background}

Browser-based molecular editors are essential components of modern
cheminformatics web applications, enabling researchers to draw, edit,
and query chemical structures without installing desktop
software~\citep{Bienfait2013,Ertl2010}.  Several capable editors
exist---JSME~\citep{Bienfait2013}, Ketcher~\citep{Ketcher2024}, and
ChemDoodle Web Components~\citep{Burger2015}---yet they share a common
architectural limitation: drawing and cheminformatics intelligence are
decoupled.  MCS computation, fingerprint generation, and stereochemistry
assignment all require a server-side toolkit or a separate WebAssembly
module, introducing latency and precluding fully offline workflows.

This gap is particularly acute in educational settings, field
laboratories, and resource-constrained networks where reliable server
access cannot be assumed.  Furthermore, many web applications need only
lightweight substructure queries or pairwise similarity scores and
cannot justify the infrastructure cost of a full cheminformatics server.

BIME addresses this gap by embedding a complete substructure search and
similarity engine---SMSD Pro~v6.6.1---directly into the editor.  SMSD
Pro is an independent JavaScript evolution of the algorithmic ideas
first described in the original Small Molecule Subgraph
Detector~\citep{Rahman2009}, substantially extended with modern
algorithms including McSplit partition refinement, VF2++ backtracking,
and multi-level coverage-driven MCS funnelling.  The result is a single
JavaScript file that delivers drawing, I/O, stereochemistry, layout, MCS,
fingerprints, and similarity search with no external dependencies.

% ====================================================================
\section{Implementation}

\subsection{Architecture}

BIME is implemented in approximately 25\,000 lines of ES5 JavaScript,
organised into 23 self-contained modules following the Immediately
Invoked Function Expression (IIFE) pattern.  Each module exposes its
public API on the \texttt{window} object.  The modules are concatenated
into a single distributable file; no transpiler, bundler, or package
manager is required.  This design ensures compatibility with all modern
browsers (Chrome, Firefox, Safari, Edge) and allows BIME to run from
\texttt{file://} URIs with no web server.

\subsection{Molecular Graph Model}

The core data structure (\texttt{Molecule.js}) represents molecules as
an atom--bond graph with 2D coordinates.  Atoms carry element symbol,
charge, isotope, chirality (\texttt{@}/\texttt{@@}), hydrogen count, and
atom-map number.  Bonds carry type (single, double, triple, aromatic),
stereo (wedge, dash), and directional markers (\texttt{/},
\texttt{\char`\\}) for E/Z specification.  Reaction SMILES are supported
natively via \texttt{>>} separation of reactant and product components.

\subsection{SMSD Pro Engine}

Six modules implement the SMSD Pro engine:

\begin{description}[nosep,leftmargin=1em]
\item[\texttt{SmsdGraph}] Converts the editor's molecular graph into
  flat adjacency arrays with atom labels, bond labels, degree sequences,
  and Neighbourhood Label Frequency (NLF) vectors for fast
  compatibility checking.
\item[\texttt{SmsdVF2}] VF2++ substructure isomorphism with FASTiso
  ordering for small molecules ($<$30 atoms) and VF3-Light ordering for
  larger targets, NLF-1/2 pruning, and a greedy $O(N)$ fast-path.
\item[\texttt{SmsdMCS}] A coverage-driven MCS funnel with seven+
  algorithmic levels: label-frequency upper bound, linear-chain and tree
  fast-paths, greedy probe, containment check, augmenting-path
  refinement, McSplit partition refinement with RRSplit maximality,
  Bron--Kerbosch with Tomita pivoting and $k$-core reduction, and
  McGregor bond-grow extension.
\item[\texttt{SmsdRings}] Smallest Set of Smallest Rings (SSSR)
  via Horton candidate generation and GF(2) Gaussian elimination,
  Relevant Cycle Basis, and Unique Ring Families.
\item[\texttt{SmsdBatch}] Batch substructure search, batch MCS, RASCAL
  fingerprint pre-screening, path-based fingerprints, ECFP/FCFP
  circular fingerprints, topological torsion fingerprints, and six
  similarity metrics (Tanimoto, Dice, cosine, Soergel, count-Tanimoto,
  count-Dice).
\item[\texttt{SmsdLayout}] SSSR-aware layout utilities including ring
  system ordering, simulated-annealing crossing reduction, PrEd-inspired
  force-directed refinement~\citep{Bertault2000}, and SMACOF stress
  majorisation~\citep{deLeeuw1977}.
\end{description}

\subsection{CIP Stereochemistry}

The \texttt{CipStereo} module implements Rules~1--4a of the IUPAC 2013
CIP recommendations~\citep{Favre2013}.  A digraph is constructed by BFS
expansion from each stereocentre with phantom-node duplication for
double/triple bonds and ring back-edges.  Priority comparison proceeds
level by level.  Rule~3 resolves ties via seqCis/seqTrans descriptors
($Z > E > $ nonstereogenic); Rule~4a applies a two-pass
like/unlike comparison for chiral centres.  Depth and node limits
(\texttt{MAX\_DEPTH}~=~30, \texttt{MAX\_NODES}~=~5000) guarantee
termination on pathological inputs.  R/S labels are assigned to
tetrahedral centres and E/Z labels to restricted double bonds.

\subsection{2D Layout}

The layout engine (\texttt{Layout.js}) generates publication-quality
coordinates following the Structure Diagram Generation (SDG)
architecture~\citep{Helson1999}: (1)~SSSR ring perception,
(2)~ring-system fusion with regular-polygon placement, edge-reflection
for fused rings, and spiro-junction handling, (3)~120-degree zigzag
chain extension, (4)~substituent fan placement, and (5)~iterative
overlap resolution using a spatial hash grid for $O(N)$ collision
detection.  Forty-five pre-computed scaffold templates
(\texttt{Templates.js}) provide instant coordinates for common ring
systems.

\subsection{SMARTS Support}

A recursive-descent SMARTS parser (\texttt{SmartsParser.js}) produces
query molecules carrying per-atom constraint arrays (wildcard, atomic
number, ring membership, degree, valence, hydrogen count, charge,
aromaticity, logical AND/OR/NOT, recursive SMARTS).  The VF2-based
matcher (\texttt{SmartsMatch.js}) evaluates these constraints during
backtracking.  Round-trip fidelity is maintained via
\texttt{SmartsWriter.js}.

% ====================================================================
\section{Features}

\subsection{Drawing and Editing}

BIME provides a full drawing toolkit: single, double, and triple bonds;
wedge/dash stereo bonds; chain drawing with angle snapping at 30-degree
intervals; ring templates (3--8-membered rings, fused bicyclics); atom
label editing with auto-valence hydrogen adjustment; charge and isotope
modification; and an unlimited undo/redo stack (100 levels).

\subsection{Chemical I/O}

Import and export are supported for SMILES (OpenSMILES specification
with full stereochemistry), reaction SMILES (\texttt{>>}), SMARTS, MOL
V2000, and MOL V3000 (automatic fallback for molecules exceeding 999
atoms).  The canonical SMILES writer uses Morgan-algorithm ordering and
H\"uckel $4n+2$ aromatic perception.

\subsection{Image Export}

Publication-quality SVG and PNG export with configurable dimensions,
padding, background colour, and font.  Clipboard copy and batch export
are supported.

\subsection{Built-in Molecule Library}

A library of 552 structures---covering common drugs, natural products,
amino acids, nucleotides, vitamins, and reagents---sourced from PubChem,
DrugBank, and ChEMBL provides instant access to frequently used
molecules.

\subsection{Feature Comparison}

Table~\ref{tab:comparison} compares BIME with three widely used
browser-based molecular editors.

\begin{table*}[t]
\centering
\caption{Feature comparison of browser-based molecular editors.
  MCS = maximum common substructure; FP = molecular fingerprints;
  CIP = Cahn--Ingold--Prelog stereo assignment.}
\label{tab:comparison}
\small
\begin{tabular}{@{}lcccc@{}}
\toprule
Feature & BIME & JSME & Ketcher & ChemDoodle Web \\
\midrule
Pure client-side (no server) & Yes & Yes & Partial\textsuperscript{a} & Yes \\
MCS (built-in) & Yes & No & No & No \\
Substructure search (built-in) & Yes & No & Server & No \\
Fingerprints (built-in) & Yes & No & No & No \\
CIP R/S and E/Z assignment & Yes & No & Server & No \\
SMARTS support & Yes & No & Yes & Yes \\
MOL V3000 & Yes & No & Yes & No \\
Reaction SMILES & Yes & Yes & Yes & Yes \\
Offline (\texttt{file://}) & Yes & Yes & No & Yes \\
Zero dependencies & Yes & Yes & No & No \\
Single-file embed & Yes & Yes & No & No \\
Open source & Apache 2.0 & BSD & Apache 2.0 & Proprietary \\
\bottomrule
\multicolumn{5}{@{}l}{\textsuperscript{a}Ketcher requires Indigo server for stereo, substructure, and layout.}
\end{tabular}
\end{table*}

% ====================================================================
\section{Results and Discussion}

\subsection{Performance}

All benchmarks were performed in Chrome 124 on a 2.6\,GHz Apple M2
MacBook Pro.  MCS computation for typical drug-like molecule pairs
(20--50 heavy atoms) completes in under 50\,ms.  The multi-level funnel
architecture ensures that simple cases (substructure containment,
identical scaffolds) are resolved at early levels without engaging the
more expensive McSplit or Bron--Kerbosch algorithms.  The adaptive
timeout (default 10\,s) guarantees responsiveness even for
worst-case inputs.

Full 2D layout generation for taxol (113 atoms, 11 rings) completes in
66\,ms, including SSSR perception, ring-system fusion, chain layout,
and overlap resolution.  Template matching bypasses the full layout
pipeline for the 45 most common scaffolds, providing sub-millisecond
coordinate generation.

\subsection{Crossing Reduction}

The simulated-annealing crossing reduction stage, combined with
PrEd-inspired force-directed refinement, substantially improves diagram
clarity.  For morphine (5 rings, 40 atoms), bond crossings are reduced
from 10 to 2.  A two-phase strategy first optimises ring-system
orientations globally, then refines individual ring flips, avoiding
local minima that trap single-phase approaches.

\subsection{CIP Correctness}

The CIP implementation was validated against the canonical examples in
the IUPAC 2013 recommendations~\citep{Favre2013} and a curated set of
challenging molecules including allenes, atropisomers (via E/Z on
restricted bonds), and spiro centres.  The digraph-based approach with
phantom-node expansion correctly handles cases where path-based
heuristics used by simpler implementations fail, such as bridged
bicyclic systems with multiple stereocentres.

\subsection{Fingerprint Quality}

ECFP and FCFP circular fingerprints follow the Morgan/extended
connectivity algorithm with configurable radius.  Path-based
fingerprints enumerate all paths up to a specified length and hash atom
and bond labels.  Tanimoto similarity scores computed by BIME on the
ChEMBL kinase inhibitor set correlate with server-computed values
(\mbox{$r^2 > 0.99$}).

\subsection{Deployment}

BIME has been deployed in production for molecular search and
educational applications.  The single-file architecture and absence of
server requirements simplify deployment to a static file host or even a
local filesystem.  Docker and shell scripts are provided for
self-hosting.

% ====================================================================
\section{Conclusions}

BIME is a lightweight, dependency-free molecular editor that brings
substructure intelligence---MCS, fingerprints, similarity search, and
CIP stereochemistry---directly to the browser.  Its single-file
deployment model makes it uniquely suited to offline, educational, and
resource-constrained settings.  BIME is available under the Apache
License 2.0 at \url{https://asad.github.io/bime-dist/}.  Future work
will extend CIP coverage to Rules~4b--5, add pharmacophore perception,
and integrate WebAssembly acceleration for molecules exceeding 200
heavy atoms.  Community contributions are welcome via the GitHub
repository.

% ====================================================================
\section*{Availability and Requirements}

\begin{description}[nosep,leftmargin=1em,font=\normalfont\bfseries]
\item[Project name:] BIME --- BioInception Molecular Editor
\item[Project home page:] \url{https://asad.github.io/bime-dist/}
\item[Operating system(s):] Platform independent (browser-based)
\item[Programming language:] JavaScript (ES5)
\item[Other requirements:] Any modern web browser (Chrome, Firefox, Safari, Edge)
\item[License:] Apache License 2.0
\item[Any restrictions to use by non-academics:] None
\end{description}

% ====================================================================
\section*{Abbreviations}

\noindent
BIME: BioInception Molecular Editor;
CIP: Cahn--Ingold--Prelog;
ECFP: Extended-Connectivity Fingerprint;
FCFP: Functional-Class Fingerprint;
IIFE: Immediately Invoked Function Expression;
MCS: Maximum Common Substructure;
NLF: Neighbourhood Label Frequency;
SDG: Structure Diagram Generation;
SMARTS: SMILES Arbitrary Target Specification;
SMILES: Simplified Molecular-Input Line-Entry System;
SMSD: Small Molecule Subgraph Detector;
SSSR: Smallest Set of Smallest Rings;
SVG: Scalable Vector Graphics;
VF2++: an improved subgraph isomorphism algorithm

% ====================================================================
\section*{Acknowledgements}

The author acknowledges the algorithmic foundations laid by the original
Small Molecule Subgraph Detector~\citep{Rahman2009}.  SMSD Pro is an
independent reimplementation and substantial extension by BioInception
PVT LTD, not a fork or derivative of the original Java codebase.

% ====================================================================
\begin{thebibliography}{10}

\bibitem[{Rahman et~al.(2009)}]{Rahman2009}
Rahman, S.A., Bashton, M., Holliday, G.L., Schrader, R., Thornton, J.M.
\newblock Small Molecule Subgraph Detector (SMSD) toolkit.
\newblock \emph{Journal of Cheminformatics}, 1:12, 2009.
\newblock \href{https://doi.org/10.1186/1758-2946-1-12}{doi:10.1186/1758-2946-1-12}

\bibitem[{Bienfait and Ertl(2013)}]{Bienfait2013}
Bienfait, B., Ertl, P.
\newblock JSME: a free molecule editor in JavaScript.
\newblock \emph{Journal of Cheminformatics}, 5:24, 2013.
\newblock \href{https://doi.org/10.1186/1758-2946-5-24}{doi:10.1186/1758-2946-5-24}

\bibitem[{Ertl(2010)}]{Ertl2010}
Ertl, P.
\newblock Molecular structure input on the web.
\newblock \emph{Journal of Cheminformatics}, 2:1, 2010.
\newblock \href{https://doi.org/10.1186/1758-2946-2-1}{doi:10.1186/1758-2946-2-1}

\bibitem[{Burger et~al.(2015)}]{Burger2015}
Burger, M.C.
\newblock ChemDoodle Web Components: HTML5 toolkit for chemical graphics,
  interfaces, and informatics.
\newblock \emph{Journal of Cheminformatics}, 7:35, 2015.
\newblock \href{https://doi.org/10.1186/s13321-015-0085-3}{doi:10.1186/s13321-015-0085-3}

\bibitem[{EPAM(2024)}]{Ketcher2024}
EPAM Systems.
\newblock Ketcher: web-based molecular editor.
\newblock \url{https://github.com/epam/ketcher}, 2024.

\bibitem[{Bertault(2000)}]{Bertault2000}
Bertault, F.
\newblock A force-directed algorithm that preserves edge-crossing properties.
\newblock \emph{Information Processing Letters}, 74(1--2):7--13, 2000.
\newblock \href{https://doi.org/10.1016/S0020-0190(00)00028-8}{doi:10.1016/S0020-0190(00)00028-8}

\bibitem[{de~Leeuw and Mair(1977)}]{deLeeuw1977}
de~Leeuw, J.
\newblock Applications of convex analysis to multidimensional scaling.
\newblock In \emph{Recent Developments in Statistics}, pages 133--145.
  North-Holland, 1977.

\bibitem[{Favre et~al.(2013)}]{Favre2013}
Favre, H.A., Powell, W.H.
\newblock \emph{Nomenclature of Organic Chemistry: IUPAC Recommendations and
  Preferred Names 2013}.
\newblock Royal Society of Chemistry, 2013.
\newblock \href{https://doi.org/10.1039/9781849733069}{doi:10.1039/9781849733069}

\bibitem[{Helson(1999)}]{Helson1999}
Helson, H.E.
\newblock Structure diagram generation.
\newblock \emph{Reviews in Computational Chemistry}, 13:313--398, 1999.
\newblock \href{https://doi.org/10.1002/9780470125908.ch6}{doi:10.1002/9780470125908.ch6}

\end{thebibliography}

\end{document}
